\documentclass[11pt,a4paper]{article}
\usepackage[T1]{fontenc}
\usepackage[utf8]{inputenc}
\usepackage{lmodern,microtype,textcomp}
\usepackage[margin=27mm,headheight=14pt]{geometry}
\usepackage{amsmath,amssymb,amsthm,mathtools}
\usepackage{booktabs,longtable,tabularx,array}
\usepackage{enumitem,listings,xcolor}
\usepackage[hidelinks,breaklinks]{hyperref}
\usepackage{fancyhdr,needspace}
\hypersetup{pdftitle={Forge Beyond Algebra: Analytic Extensions},pdfauthor={Research proposal prepared for Vladimir Reshetnikov},pdfsubject={Anchored analytic certificates, differential ladders, and rational series tails}}
\pagestyle{fancy}\fancyhf{}
\fancyhead[L]{\small Forge: analytic extensions}
\fancyhead[R]{\small Research proposal}
\fancyfoot[C]{\thepage}
\setlength{\parskip}{0.32em}
\setlength{\parindent}{1em}
\setlength{\emergencystretch}{2em}
\setlist{nosep,leftmargin=1.6em}
\lstset{basicstyle=\ttfamily\footnotesize,breaklines=true,columns=fullflexible,
 keepspaces=true,showstringspaces=false,upquote=true,frame=single,framerule=.25pt,
 xleftmargin=.3em,xrightmargin=.3em,aboveskip=.7em,belowskip=.7em}
\newtheorem{theorem}{Theorem}[section]
\newtheorem{proposition}[theorem]{Proposition}
\newtheorem{lemma}[theorem]{Lemma}
\newtheorem{corollary}[theorem]{Corollary}
\theoremstyle{definition}\newtheorem{definition}[theorem]{Definition}
\newtheorem{example}[theorem]{Example}
\newtheorem{remark}[theorem]{Remark}
\newcolumntype{P}[1]{>{\raggedright\arraybackslash}p{#1}}
\newcommand{\R}{\mathbb R}\newcommand{\Q}{\mathbb Q}\newcommand{\N}{\mathbb N}
\newcommand{\code}[1]{\texttt{\detokenize{#1}}}
\newcommand{\forge}{\texttt{forge}}
\newcommand{\dd}{\mathrm d}
\newcommand{\I}{\mathcal I}
\newcommand{\Eval}{\operatorname{eval}}
\newcommand{\Unknown}{\textsf{unknown}}
\newcommand{\status}[1]{\par\noindent\textbf{Evidence status.} #1\par}
\begin{document}
\begin{titlepage}
\vspace*{1.2cm}
\noindent{\Large\bfseries Forge Beyond Algebra}\par
\vspace{.55cm}
\noindent{\LARGE Anchored analytic certificates, integrating-factor search,\par
and quantified infinite tails}\par
\vspace{1cm}
\noindent{\large Concrete extensions to the audited Forge design}\par
\vspace{.7cm}
\noindent Prepared for Vladimir Reshetnikov\par
\noindent 15 September 2026\par
\vspace{.8cm}
\begin{abstract}
The current Forge design already supplies a substantial collection of structural,
algebraic, finite-state, and finite-summation workers. Repeating those proposals
would add little. This article develops three complementary analytic workers:
exact anchored remainders that preserve a function's order of vanishing;
search for integrating-factor differential ladders proving inequalities on an
unbounded half-line; and rational tail barriers that establish summability,
explicit tail estimates, and executable Cauchy moduli. A companion harmonic
minorant supplies a genuinely different negative outcome: certified divergence
of eventually positive rational series.

The proposal specifies certificate languages, search algorithms, mathematical
soundness proofs, limitations, library interfaces, and an integration sequence.
For the rational-series fragment, an elementary positive-shift argument gives a
constructive convergence/divergence classification and an explicit certificate
construction, rather than a bounded search heuristic alone. Existing Taylor-model
and verified-numerics work is credited and reused; these are extensions relative
to the inspected Forge snapshot, not claims to have invented interval arithmetic,
Taylor models, integrating factors, or convergence comparison.

The accompanying standard-library Python implementation produced 246 accepted
certificate records, representing 230 distinct source/domain subjects. A separate
sparse-polynomial checker replays them with search imports blocked; 738
specified mutations are rejected. Symbolic cross-checks, numerical enclosure
diagnostics, and exact least-index witness checks accompany the corpus. These
are component experiments, not Lean proofs or tactic benchmarks. No Lean
executable was available in the authoring environment, so the Lean bridge
specimens and all end-to-end integrations remain explicitly uncompiled.
\end{abstract}
\vfill
\noindent\small Audited Forge revision:\par
\noindent\path{c98e47c5f804e92880fc1d0e378c1b95832b685c}\par
\medskip
\noindent\small The archive includes editable \LaTeX, this PDF, executable searches,
standalone replay, tests, certificates, measurements, and implementation notes.
\end{titlepage}
\tableofcontents
\newpage

\section{The contribution, and what is deliberately not repeated}
\label{sec:delta}

The reviewed Forge repository describes a certificate-producing proof-planning
layer above Lean's \code{grind}, not an installed tactic. Its merged article
incorporates eighteen proposal packages. The extension round already includes
observable-space closure, polynomial-ideal closure, finite-algebra covers,
finite Kripke countermodels, boundary-safe telescoping, and Ore-operator
transport. The earlier round already includes induction and generalisation,
witness synthesis, nonlinear polynomial certificates, and Boolean/theory
cooperation~\cite{forge-readme,forge-ledger,forge-conclusion}.

The right question is therefore not ``what other solver can we name?'' It is:
\emph{which important proof obligations still require a mathematical representation
that the audited workers do not provide?} The proposal here addresses analytic
contact, global differential comparison, and infinite convergence. The common
pattern is to turn an infinite family of analytic obligations into a finite
algebraic certificate, while retaining the theorem that licenses that reduction.

\subsection{A repository-relative delta ledger}

\begin{longtable}{@{}P{.25\textwidth}P{.31\textwidth}P{.37\textwidth}@{}}
\toprule
Already present or established & Not claimed as a contribution & Addition developed here \\
\midrule\endhead
Polynomial cones, quadratic certificates, Bernstein subdivision, square-factor/Sturm reasoning &
Another polynomial positivity engine &
A source-preserving analytic remainder that exposes a \emph{polynomial quotient} without losing the zero at the anchor. \\
Invariant ideals and finite-dimensional closures &
Renaming invariant generation &
A first-order differential comparison certificate that reconstructs positivity on all of $[0,\infty)$ from finitely many exact identities. \\
Finite Gosper/creative telescoping and boundary-safe sums &
Another finite summation identity &
A tail potential proving summability before using an infinite sum, together with a strict Cauchy modulus and a divergence alternative. \\
Arithmetic witness synthesis &
The general idea of synthesising witnesses &
An explicit least-index construction for a certified tail bound, checked by two rational inequalities. \\
Source-owned evidence, opaque receipts, scope discipline, and re-elaboration &
Another general trust architecture &
The new source obligations: analytic domain, anchor, interval orientation, natural-index conversion, eventual sequence identity, and summability. \\
Verified numerics and Taylor models in the wider literature &
Inventing interval arithmetic or Taylor arithmetic &
A Forge-specific routing and certificate layer, with concrete reuse decisions and new fragment-specific algorithms. \\
\bottomrule
\end{longtable}

The audit used the repository README, the merged baseline section, the
provenance ledger, the conclusion, the relevant nonlinear section, and the
Leant/Djex review. It is \emph{not} a claim that every line of all eighteen
original submissions was reread. The absence claim is correspondingly scoped:
these three workers were not found as developed capabilities in the audited
merged material~\cite{forge-baseline,forge-nonlinear,forge-neighbours}.

\subsection{Why these are useful proof shapes}

An algebraic worker may regard $\sin x$ as an opaque real atom. That is safe,
but it does not know that $x-\sin x$ has a cubic zero at the origin. A bounded
numerical worker may prove an inequality on $[0,1]$ while saying nothing about
$[1,\infty)$. A finite telescoping identity can be valid even when an infinite
series diverges. Each gap requires a different mathematical bridge, not a
larger search budget.

The proposed output is stronger than a yes/no heuristic. An anchored worker can
return $f(x)\ge\delta x^m$ on an exact interval. A ladder worker can return a
chain of differential identities proving a global sign. A tail worker can return
an explicit $N$, a bound on every remaining finite partial sum, summability,
and an executable modulus. Its divergence outcome carries a harmonic comparison,
not merely an exhausted convergence search.

\subsection{What the experiments establish}

The implemented searches and replay checker use exact rational arithmetic. The
search polynomial representation is dense; the checker has a separate sparse
implementation and imports no search module. The corpus has 62 anchored records,
43 ladder records, 81 barrier records, and 60 divergence records. Two elementary
examples are repeated across named subfamilies, and seven series are optimised
at three different cutoffs; consequently there are 230 distinct source/domain
subjects, not 246 distinct theorems. Section~\ref{sec:experiments} makes the units
and these overlaps explicit.

No result here establishes a speedup over \code{grind}, Mathlib tactics, or
LeanCert. Small examples may already have a library lemma or an easy manual
proof. Their purpose is to exhibit the certificate and the search, not to make
an existing tactic look weak. Nor does another collection of Python experiments
close Forge's existing Lean integration gap. The repository's own conclusion
correctly distinguishes algorithm design from shipping a checked end-to-end
pipeline~\cite{forge-conclusion}.

\section{A worked example: do not erase the zero}
\label{sec:example}

Consider proving $x-\sin x\ge0$ for $0\le x\le1/2$. A fourth-order model is
\[
 \sin x=x-\frac{x^3}{6}+x^4r(x),\qquad |r(x)|\le\frac5{108}.
\]
The rational remainder bound follows from the majorant in
Section~\ref{sec:primitives}. Subtracting gives
\[
 x-\sin x=x^3\left(\frac16-xr(x)\right).
\]
Since $x\ge0$ and $x\le1/2$,
\[
 \frac16-xr(x)\ge\frac16-\frac12\frac5{108}=\frac{31}{216}>0.
\]
Thus the same certificate proves the stronger fact
\begin{equation}\label{eq:sine-output}
 x-\sin x\ge\frac{31}{216}x^3\qquad(0\le x\le1/2).
\end{equation}
In particular it also proves strict positivity whenever $x>0$.

The difference between this and an unstructured interval calculation is exact.
If we replace $x^4r(x)$ by a uniform error interval \emph{before} factoring, the
lower bound becomes
\[
 0-\left(\frac12\right)^4\frac5{108}=-\frac5{1728}<0.
\]
No contradiction has occurred: the lower bound is valid but too weak. It has
forgotten that both the main term and the error vanish at zero at different
rates. Preserving that information changes what can be certified without
excluding the endpoint.

There is no division by $x$ in the certificate theorem. The factorisation is an
identity, and $x=0$ is covered directly. This is important in Lean: replacing
$f(x)$ by $f(x)/x^3$ and forgetting a nonzero side condition would prove a
different statement. The representation below avoids that trap structurally.

This example is a demonstration of the worker, not a claim that the basic
inequality for sine is missing from Mathlib. The general utility appears when
several analytic subexpressions cancel to high order and the appropriate
quotient is not a named theorem.

\section{Worker A: exact anchored remainder arithmetic}
\label{sec:anchor}
\status{Implemented in \code{jets.py}; independently replayed in Python.
The semantic theorem is proved below on paper, not in Lean.}

\subsection{The semantic object}

\begin{definition}[Anchored enclosure]
Fix $b>0$ and $K\ge1$. An anchored enclosure of $f:[0,b]\to\R$ is a pair
$(p,I)$, where $p\in\Q[x]$ has degree less than $K$ and $I=[\ell,u]$ is a
rational interval, such that
\begin{equation}\label{eq:anchorsem}
 \forall x\in[0,b],\quad \exists r\in[\ell,u],\quad
 f(x)=p(x)+x^Kr.
\end{equation}
\end{definition}

The existential formulation deliberately does not require dividing by $x^K$.
It also does not require a globally selected remainder function. At zero, only
the exact constant coefficient matters. Away from zero, ordinary analytic
remainder formulas supply the witness.

This is a standard Taylor-remainder representation, not a new numerical
abstraction. Streeter and Dillon explicitly study intervals multiplying
$(x-x_0)^K$ and compositional Taylor-mode bounds. Earlier verified interval
work also combines elementary functions with Taylor estimates
~\cite{autobound,verified-intervals}. Our contribution is its certificate use
and goal-directed order-of-vanishing extraction in Forge.

Let $H_b(p)$ denote ordinary rational Horner interval evaluation on $[0,b]$.
The essential property is
\[
 p([0,b])\subseteq H_b(p).
\]
All interval endpoints in the prototype are exact fractions. Floating-point
arithmetic does not participate in search acceptance or replay.

\subsection{Primitive certificates}
\label{sec:primitives}

The prototype supports rational polynomials and the primitives
\[
 e^{ax},\quad\sin(ax),\quad\cos(ax),\quad\log(1+ax),\quad\arctan(ax),
 \qquad a\in\Q,
\]
combined by addition, negation, and multiplication. A node labelled
\code{log1p} has the exact meaning $\log(1+ax)$; it is not an arbitrary logarithm
node whose domain is guessed later.

For exponential, sine, and cosine, use the ordinary degree-$(K-1)$ Taylor
polynomial and the symmetric remainder radius
\begin{equation}\label{eq:expR}
 R_K(a,b)=\frac{|a|^K}{K!\bigl(1-|a|b/(K+1)\bigr)},
 \qquad \frac{|a|b}{K+1}<1.
\end{equation}
For logarithm and arctangent, use
\begin{equation}\label{eq:logR}
 S_K(a,b)=\frac{|a|^K}{K(1-|a|b)},\qquad |a|b<1.
\end{equation}
The latter guard also ensures $1+ax>0$ for the logarithm. The radius for
arctangent is deliberately conservative: it ignores parity and alternating
signs in exchange for a tiny shared checking rule.

\begin{lemma}[Primitive enclosures]
With the corresponding guard, the indicated Taylor polynomial and interval
$[-R_K,R_K]$ or $[-S_K,S_K]$ satisfy \eqref{eq:anchorsem}.
\end{lemma}
\begin{proof}
For the exponential tail and $0<x\le b$, division by $x^K$ \emph{inside this
mathematical proof} gives the absolute majorant
\[
 \sum_{j=K}^{\infty}\frac{|a|^jx^{j-K}}{j!}.
\]
Its first term is $|a|^K/K!$. Each successive ratio is at most
$|a|b/(K+1)<1$, giving \eqref{eq:expR}. The sine and cosine tails are
subseries of the same absolute majorant, with signs discarded.

For logarithm or arctangent, the absolute coefficient in degree $j\ge K$ is
at most $|a|^j/j\le |a|^j/K$. Summing the resulting geometric majorant after
factoring $x^K$ gives \eqref{eq:logR}. The power series represent the stated
functions under the guards. At $x=0$ the truncated polynomial has the correct
value, so choosing any remainder in the nonempty interval completes the proof
without division at zero.
\end{proof}

A production library should replace these symmetric bounds by stronger signed
or minimax bounds where available. Nothing in the certificate interface requires
this particular approximation. The current implementation makes an intentionally
small, auditable choice; it does not compete with a mature numerical library's
full catalogue of primitive bounds.

\subsection{Composition rules}

Addition and negation act componentwise. For a polynomial input $s$, split
$s=p+x^Kh$ and use $H_b(h)$ as the remainder interval. The product rule is the
only nontrivial compositional case.

Suppose $f=p+x^Kr$ and $g=q+x^Ks$, with $r\in I$ and $s\in J$. Write
\[
 pq=t+x^Kh,\qquad \deg t<K.
\]
Then
\begin{equation}\label{eq:jetproduct}
 fg=t+x^K\bigl(h+ps+qr+x^Krs\bigr).
\end{equation}
Consequently a valid product remainder interval is
\begin{equation}\label{eq:productinterval}
 H_b(h)+H_b(p)J+H_b(q)I+[0,b^K]IJ.
\end{equation}
All operations in this expression are interval operations. They preserve
containment even when occurrences are correlated; the cost of forgetting a
correlation is loss of precision, not unsoundness.

\begin{theorem}[Compositional soundness]
For every accepted expression in the prototype grammar, the enclosure computed
by these rules satisfies \eqref{eq:anchorsem}.
\end{theorem}
\begin{proof}
Induct on the expression tree. Polynomial nodes follow from exact coefficient
splitting and Horner containment. Primitive nodes follow from the preceding
lemma. Addition and negation follow by combining the existential remainder
witnesses at the same $x$. For multiplication use \eqref{eq:jetproduct} and
containment in each of the four intervals in \eqref{eq:productinterval}.
\end{proof}

\subsection{Extracting the order of vanishing}

Let $m<K$ be the least index of a nonzero coefficient of $p$. Then
$p=x^mq$ and
\[
 f(x)=x^m\bigl(q(x)+x^{K-m}r\bigr).
\]
Compute
\begin{equation}\label{eq:margin}
 \delta=\inf\bigl(H_b(q)+[0,b^{K-m}]I\bigr).
\end{equation}
If $p=0$, instead set $m=K$ and $\delta=\inf I$.

\begin{theorem}[Anchored positivity certificate]\label{thm:anchored}
An accepted enclosure with $\delta\ge0$ proves
\[
 f(x)\ge\delta x^m\ge0\qquad\text{for every }x\in[0,b].
\]
\end{theorem}
\begin{proof}
When $p\ne0$, \eqref{eq:margin} is a lower bound on the factor in parentheses.
Multiplication by $x^m\ge0$ preserves the inequality. When $p=0$, the conclusion
follows directly from $f(x)=x^Kr$ and $r\ge\delta$. Both arguments include zero.
\end{proof}

The certificate stores the expression, exact bound $b$, order $K$, polynomial,
remainder interval, valuation $m$, and margin $\delta$. The checker recomputes
all of them from the externally supplied subject. Neither a claimed zero
coefficient nor a claimed margin is trusted.

\subsection{Search, cost, and a useful completeness statement}

The implemented search tries $K=2,3,\ldots,K_{\max}$ at the supplied $b$ and
returns the first nonnegative margin. Unsupported primitive guards are skipped.
The checker admits at most 256 expression nodes, depth 32, and order 40. Dense
truncated polynomial multiplication costs $O(K^2)$ rational operations per
multiplication node. Exact numerator and denominator bit growth must be
accounted for separately from operation counts.

\begin{lstlisting}
anchored_search(expression, exact_box, order_budget):
    for K in 2 .. order_budget:
        model := build_anchored_model(expression, exact_box, K)
        if a primitive guard is unavailable: continue
        m := first nonzero polynomial coefficient, or K
        margin := lower bound after extracting x^m
        if margin >= 0:
            return candidate(model, m, margin)
    return unknown
\end{lstlisting}

For a goal asking for a neighbourhood rather than a prescribed interval, the
controller may additionally search $b_0/2^j$. This wrapper is specified, not
separately executed in the corpus.

\begin{proposition}[Local positive-germ completeness, without resource caps]
For the supported analytic grammar, suppose the first nonzero Taylor
coefficient of $f$ at zero has positive value and degree $m$. A fair search over
orders and shrinking positive dyadic radii eventually obtains an anchored
positivity certificate.
\end{proposition}
\begin{proof}
Choose a fixed order $K>m$. On sufficiently small radii all primitive guards
hold and all propagated remainder intervals remain uniformly bounded. This
follows inductively from the explicit primitive radii and the algebraic product
rule. The Horner interval for $q=p/x^m$ tends to the positive constant $q(0)$
as $b\to0$. Meanwhile $[0,b^{K-m}]I$ tends to $\{0\}$. Therefore the lower
endpoint in \eqref{eq:margin} is eventually positive. A fair search reaches such
an order/radius pair.
\end{proof}

This theorem does not say that the fixed-radius implementation decides analytic
inequalities. Nor does it decide equality of analytic functions. For example,
$\sin^2x+\cos^2x-1=0$ is true, but the prototype's independently propagated
symmetric errors do not cancel exactly; its nonnegativity search returns
\Unknown{} under the recorded cap. The existing algebraic/rewrite layer should
recognise the identity before numerical approximation is attempted.

\section{Worker B: integrating-factor differential ladders}
\label{sec:ladder}
\status{Exact search and Python replay are implemented in \code{ladders.py}
and \code{checker.py}; derivative identities are independently cross-checked
with SymPy. Lean reconstruction is specified, not executed.}

\subsection{A global proof rule}

Let $D$ denote differentiation and define
\[
 L_\alpha=D-\alpha,\qquad L_\alpha f=f'-\alpha f.
\]

\begin{lemma}[One integrating-factor step]\label{lem:factorstep}
Let $f$ be continuously differentiable on $[0,\infty)$, let $\alpha\in\R$,
and suppose $f(0)\ge0$ and $L_\alpha f(x)\ge0$ for all $x\ge0$. Then
$f(x)\ge0$ for all $x\ge0$.
\end{lemma}
\begin{proof}
For $h(x)=e^{-\alpha x}f(x)$ the product rule gives
\[
 h'(x)=e^{-\alpha x}\bigl(f'(x)-\alpha f(x)\bigr)\ge0.
\]
Thus $h(x)\ge h(0)=f(0)\ge0$. Multiply by the positive number $e^{\alpha x}$.
Alternatively, apply the fundamental theorem of calculus on each finite
interval $[0,x]$; no compactness assertion about the entire half-line is needed.
\end{proof}

An \emph{integrating-factor ladder} is a finite sequence
\[
 f_0=f,\qquad f_{i+1}=L_{\alpha_i}f_i\quad(0\le i<r),
\]
with $f_i(0)\ge0$ and an independently checkable proof of $f_r\ge0$ on the
half-line. Applying the lemma backwards proves the original goal.

\subsection{An exact finite language}

The implemented language is the rational exponential-polynomial space
\[
 f(x)=\sum_{\lambda\in\Lambda}p_\lambda(x)e^{\lambda x},
 \qquad\Lambda\subset\Q\text{ finite},\quad p_\lambda\in\Q[x].
\]
The transformation is exact and remains in the language:
\begin{equation}\label{eq:rateupdate}
 L_\alpha\bigl(p_\lambda e^{\lambda x}\bigr)
 =\bigl(p_\lambda'+(\lambda-\alpha)p_\lambda\bigr)e^{\lambda x}.
\end{equation}
A cheap terminal cone requires every coefficient of every $p_\lambda$ to be
nonnegative. Then each polynomial is nonnegative for $x\ge0$ and every
exponential is positive. This terminal cone is intentionally narrow; Forge's
existing polynomial worker can later replace it with stronger certificates.

\begin{theorem}[Ladder certificate soundness]
Suppose the certificate checker verifies every identity
\eqref{eq:rateupdate}, every boundary seed, and membership of $f_r$ in the
terminal cone. Then $f_0(x)\ge0$ for all $x\ge0$.
\end{theorem}
\begin{proof}
Every exponential polynomial is continuously differentiable. The terminal
coefficient test proves $f_r\ge0$. Apply Lemma~\ref{lem:factorstep} for
$i=r-1,r-2,\ldots,0$ using the corresponding seed and identity.
\end{proof}

\begin{example}[All orders of the exponential remainder]
For an integer $s\ge1$, set
\[
 f_s(x)=e^x-\sum_{j=0}^{s-1}\frac{x^j}{j!}.
\]
One exact transformation gives
\[
 (D-1)f_s(x)=\frac{x^{s-1}}{(s-1)!}.
\]
Since $f_s(0)=0$, a single ladder step proves $f_s(x)\ge0$ for every $x\ge0$.
The complexity of this certificate does not depend on the width of a numerical
box. For $s=2$, the stored certificate uses rate $1$ and terminal polynomial $x$.
\end{example}

\subsection{Bounded search with a small candidate set}

The search uses breadth-first exploration, offering rates in
$\Lambda\cup\{0\}$. Each state is an exact canonical exponential polynomial.
A state with a negative seed is discarded from this particular ladder search,
not declared to make the original goal false. Canonical states are memoised;
with unit-cost transitions, first discovery by breadth-first search has maximal
remaining depth among paths to that state.

\begin{lstlisting}
queue := [(source, empty_rate_list)]
while queue nonempty and node_budget remains:
    f, path := pop_front(queue)
    if coefficients of f are all nonnegative:
        return the exact stages, rates, and boundary seeds
    if depth budget exhausted: continue
    for alpha in source_rates union {0}:
        g := exact_derivative(f) - alpha * f
        if g(0) >= 0 and g is new:
            push_back(queue, (g, path + [alpha]))
return unknown
\end{lstlisting}

For $s$ rate blocks of maximum degree $d$, one transformation uses $O(sd)$
rational arithmetic operations. The naive search tree has branching factor at
most $s+1$ and can be exponential in the depth cap. The experiments use a depth
cap of 12 and a node cap of 4,000. The certificate checker accepts any admitted
rational rates, not only the heuristic candidate set used by search. Therefore
improving rate selection does not require enlarging the checker.

A concrete next rate heuristic is to inspect the leading coefficient in each
block. Choosing $\alpha=\lambda$ removes the degree-preserving term from that
block, leaving a derivative and reducing its degree. Candidate rates outside
$\Lambda$ could also be proposed by solving the linear inequalities in
$\lambda-\alpha$ required by a terminal coefficient test. That extension is
not implemented here.

\subsection{A precise incompleteness obstruction}

Consider the globally nonnegative function
\[
 F(x)=(e^x-2)^2=e^{2x}-4e^x+4.
\]
At $a=\log2>0$, $F(a)=F'(a)=0$ and $F''(a)>0$. For any finite real
$\alpha$, let $G=F'-\alpha F$. Then $G(a)=0$ and $G'(a)=F''(a)>0$.
Hence $G(x)<0$ immediately to the left of $a$. No first integrating-factor
step can have a globally nonnegative residual. A zero-step certificate also
fails the prototype's coefficient cone because the $e^x$ coefficient is $-4$.

Thus no choice of a longer ladder can repair this particular failure within the
specified proof rule: the first required implication is already obstructed.
The correct response is to expose the square and use the existing algebraic
positivity machinery. The recorded depth-eight search returns \Unknown{}, as
it should. This example is also a test against an implementation that checks
only the seeds but forgets that every residual must ultimately be proved
nonnegative on the whole domain.

\subsection{What is and is not new}

The mathematical rule is ordinary integrating-factor comparison. Mathlib
already provides derivative-based comparison infrastructure
~\cite{mathlib-meanvalue}. The proposed capability is automatic \emph{discovery}
of a finite chain in an exact language, followed by a small compositional
certificate. It is not a new differential-equation solver, a general
transcendental decision procedure, or a claim that exponential-polynomial
positivity is decided by this search.

\section{Worker C: rational tail barriers}
\label{sec:tails}
\status{Constructive synthesis, coefficient optimisation, Python replay,
finite-partial-sum checks, and modulus witnesses are implemented. The infinite
semantic bridge is proved on paper and remains to be formalised in Lean.}

\subsection{Why a finite telescoper is not enough}

An infinite-series theorem has obligations absent from a finite sum: the terms
must be defined on the tail, finite initial exceptions must be separated,
convergence must be established, and only then may the infinite-sum value be
used. The prototype's subject therefore says \emph{eventual rational tail}.
It does not claim to interpret an arbitrary source function at poles of a
rational expression.

Let a real sequence satisfy
\begin{equation}\label{eq:rationalsequence}
 a_n=\frac{A(n)}{D(n)}\qquad(n\ge N),
\end{equation}
where $A,D\in\Q[n]$, and let
\begin{equation}\label{eq:barrier}
 B(n)=\frac{C}{(n+c)^p},\qquad C>0,\quad p\in\N,\quad p\ge1.
\end{equation}
The source bridge must prove \eqref{eq:rationalsequence}, including the
interpretation of the natural-number index as a real number. The certificate
will prove $D(n)>0$ and $n+c\ge1$ on the admitted tail.

\begin{definition}[Tail barrier]
A tail barrier at $N$ consists of $a_n\ge0$, $B(n)\ge0$, and
\begin{equation}\label{eq:stepbarrier}
 a_n+B(n+1)\le B(n)\qquad(n\ge N).
\end{equation}
\end{definition}

\begin{theorem}[Finite and infinite consequences]\label{thm:tailsound}
Under these assumptions, for $M\ge N$,
\begin{equation}\label{eq:finitetail}
 0\le\sum_{n=N}^{M}a_n\le B(N)-B(M+1)\le B(N).
\end{equation}
The nonnegative tail series is summable, and
\[
 0\le\sum_{n=N}^{\infty}a_n\le B(N).
\]
The same statements hold with any larger starting index in place of $N$.
\end{theorem}
\begin{proof}
Sum \eqref{eq:stepbarrier} over the finite range and cancel the adjacent
$B$ terms. Nonnegativity of the $a_n$ gives the left inequality; nonnegativity
of $B(M+1)$ gives the last. Thus the tail partial sums form a nondecreasing
bounded sequence of real numbers and converge. Taking the limit in the finite
bound gives the tail estimate. Applying the same argument starting at a larger
index proves the final statement.
\end{proof}

A stronger bound on an individual finite range does not depend on any
infinite-sum convention. In Lean, it should be the first lemma proved and tested.
Summability is a separate theorem derived from it. Existing Mathlib infinite-sum
and sum--integral comparison APIs are appropriate reusable infrastructure,
not functionality to implement from scratch~\cite{mathlib-infinitesum,mathlib-sumintegral}.

\subsection{A polynomial residual with all signs accounted for}

Set $u(n)=(n+c)^p$ and $v(n)=(n+c+1)^p$. On the admitted tail, all of $D,u,v$
are positive. Multiplication by $Duv$ therefore preserves inequalities, and
\eqref{eq:stepbarrier} is equivalent to
\begin{equation}\label{eq:residual}
 R(n)=C\bigl(v(n)-u(n)\bigr)D(n)-A(n)u(n)v(n)\ge0.
\end{equation}
This orientation is easy to reverse accidentally: the correct telescoping
quantity is $B(n)-B(n+1)$.

The checker uses the shift $n=N+t$. It requires all coefficients of
$A(N+t)$, $D(N+t)$, and $R(N+t)$ to be nonnegative, with the constant
coefficient of $D(N+t)$ strictly positive. Then for $t\ge0$, numerator and
residual are nonnegative and the denominator is strictly positive. It also
checks $N+c\ge1$. These finite rational checks imply every side condition in
Theorem~\ref{thm:tailsound}.

\begin{example}[An exact telescoper]
For $a_n=1/(n(n+1))$, choose $N=1$, $C=1$, $c=0$, and $p=1$. Then
$R=0$ identically, so the checker recognises an exact telescoping identity as a
special case of a barrier. It does not need a different infinite-series rule.
\end{example}

\begin{example}[A searched inequality, not an exact telescoper]
For $a_n=1/n^2$ and $N=2$, the coefficient optimiser chooses
\[
 B(n)=\frac{3}{2n},\qquad \sum_{n=2}^{\infty}\frac1{n^2}\le\frac34.
\]
Here $R(n)=\tfrac12n^2-n$, so $R(2+t)=t+\tfrac12t^2$ has nonnegative
coefficients. This is an independently checkable rational inequality; a decimal
approximation to the sum is unnecessary.
\end{example}

\subsection{Exact optimisation of the barrier constant}

For fixed $N,c,p$, write
\[
 R(N+t)=C U(t)-V(t).
\]
For each coefficient, the requirement $C U_i-V_i\ge0$ gives a half-line or an
interval constraint on $C$. If $U_i>0$, require $C\ge V_i/U_i$; if $U_i<0$,
require $C\le V_i/U_i$; if $U_i=0$, reject the template when $V_i>0$.
Intersect all constraints with $C>0$. This solves the coefficient feasibility
problem exactly with rational comparisons and does not need a linear-programming
package.

The implementation enumerates a finite set of shifts and exponents, solves this
one-dimensional feasibility problem for each, and keeps the smallest certified
$B(N)$. Its default shift set is $\{-1,0,1\}$ and its exponent cap is eight.
It is complete for each fixed template's \emph{coefficientwise criterion}, not
for all valid rational barriers or all nonnegative polynomials. A polynomial
can be positive on a half-line despite having a negative shifted coefficient.
A stronger existing Forge polynomial certificate can replace this criterion
without changing the infinite-sum theorem.

If the feasible interval has infimum zero but no positive minimum, the
implementation chooses a positive feasible value rather than asserting that
an unattained optimum was attained. For a nonzero nonnegative tail, positive
lower constraints normally force a positive minimum; the boundary convention
still needs to be explicit.

\section{A constructive classification for rational series}
\label{sec:classification}

A bounded template search is useful, but this fragment permits more. The next
result is a \emph{certificate construction}: it explains why increasing the
cutoff eventually succeeds and gives a computable cutoff, rather than relying
on repeated trial and error.

\subsection{Making every shifted coefficient positive}

\begin{lemma}[Positive-shift lemma]\label{lem:shift}
Let $q(n)=\sum_{j=0}^{m}q_jn^j\in\Q[n]$ be nonzero with $q_m>0$.
Define, for $m>0$,
\[
 H=\max_{0\le k\le j<m}
 \frac{|\binom{j}{k}q_j|}{\binom{m}{k}q_m},
\]
and set $H=0$ for $m=0$. Then every coefficient of $q(N+t)$ is strictly
positive whenever $N\ge\lfloor H\rfloor+3$.
\end{lemma}
\begin{proof}
The coefficient of $t^k$ is
\[
 g_k(N)=\sum_{j=k}^{m}\binom{j}{k}q_jN^{j-k}.
\]
For $k=m$ this is $q_m>0$. For $k<m$, put
$h=\binom{m}{k}q_m>0$ and $d=m-k\ge1$. By the definition of $H$,
\[
 g_k(N)\ge h\left(N^d-H\sum_{i=0}^{d-1}N^i\right).
\]
Since $N>H+1$ and $N>1$,
\[
 H\sum_{i=0}^{d-1}N^i
 =H\frac{N^d-1}{N-1}<N^d.
\]
Thus $g_k(N)>0$. The constant-polynomial case is immediate.
\end{proof}

The bound is conservative. The search may find a much smaller usable cutoff,
but replay does not need to trust any root bound: it directly recomputes and
checks the final shifted coefficients. The implementation tests this helper
on 200 additional randomly generated positive-leading polynomials.

\subsection{Convergence certificates}

\begin{theorem}[Constructive upper barrier]\label{thm:constructive}
Let $A,D\in\Q[n]$ be nonzero with positive leading coefficients $a,d_0$,
and let
\[
 d=\deg D-\deg A\ge2,\qquad p=d-1.
\]
For every rational $C>a/(p d_0)$, there is an explicitly computable integer
$N\ge1$ for which the coefficient checker accepts
$B(n)=C/n^p$ as a tail barrier for $A(n)/D(n)$.
\end{theorem}
\begin{proof}
With $c=0$, the polynomial
\[
 R(n)=C\bigl((n+1)^p-n^p\bigr)D(n)-A(n)n^p(n+1)^p
\]
has degree $\deg D+p-1=\deg A+2p$. Its leading coefficient is
$C p d_0-a>0$. The polynomials $A,D,R$ therefore all have positive leading
coefficients. Apply Lemma~\ref{lem:shift} to each and choose $N$ at least the
maximum of the three resulting cutoffs and 1. Their shifted coefficients are
strictly positive, hence all barrier checks hold.
\end{proof}

The prototype chooses $C=a/(p d_0)+1$. A smaller rational slack can improve the
asymptotic constant at the expense of a potentially larger cutoff and larger
coefficient sizes. This is a genuine tradeoff between the strength of the
analytic result and the cost of its certificate.

\subsection{A certified divergence outcome}

\begin{theorem}[Constructive harmonic minorant]\label{thm:diverge}
Under the same positive-leading assumptions, suppose
$d=\deg D-\deg A\le1$. There are explicitly computable rational $c_0>0$
and integer $N\ge1$ such that
\[
 \frac{A(n)}{D(n)}\ge\frac{c_0}{n}\qquad(n\ge N).
\]
Consequently the positive tail partial sums are unbounded and the series is not
summable.
\end{theorem}
\begin{proof}
If $d=1$, choose $c_0=a/(2d_0)$. Then $nA(n)-c_0D(n)$ has positive
leading coefficient $a/2$. If $d\le0$, choose $c_0=1$; the term $nA(n)$
has degree strictly greater than $D(n)$ and positive leading coefficient.
In either case apply Lemma~\ref{lem:shift} to $D$ and
$H(n)=nA(n)-c_0D(n)$. Beyond the resulting cutoff, $D(n)>0$ and $H(n)\ge0$.
Division by the positive number $nD(n)$ yields the claimed minorant.

For an explicit divergence argument, group the harmonic tail into blocks
$2^jN\le n<2^{j+1}N$. Each block contains $2^jN$ terms, each at least
$1/(2^{j+1}N)$, so its sum is at least $1/2$. Every corresponding block of
our series contributes at least $c_0/2$. Arbitrarily many blocks give
unbounded partial sums.
\end{proof}

The checker stores $c_0$, $N$, the shifted denominator, and the shifted
polynomial $H$. It recomputes these polynomials from the external source
subject. The verdict means \emph{divergence of the stated eventual rational
tail under its bridge assumptions}. It does not mean that a failed barrier
search has disproved convergence, and it does not misuse a real-valued
\code{tsum} as though it were an extended-real infinity.

\begin{corollary}[Unbounded-algorithm completeness on the stated fragment]
For nonzero rational $A,D$ with positive leading coefficients, the construction
produces a summability certificate when $\deg D-\deg A\ge2$ and a divergence
certificate when $\deg D-\deg A\le1$. Finite alterations of source values do
not change the classification.
\end{corollary}
\begin{proof}
Apply Theorems~\ref{thm:constructive} and~\ref{thm:diverge}. A finite prefix
has a finite real sum and cannot change convergence or unbounded tail growth.
\end{proof}

The familiar convergence classification is not presented as new mathematics.
The useful algorithmic result is the explicit compilation into one of two
small certificate languages, using only rational polynomial arithmetic and a
computed shift. The prototype has resource bounds, so the corollary is a
statement about the uncapped construction, not a guarantee that every input is
accepted by the shipped bounded checker.

The zero numerator is a trivial separate case and is excluded by the current
constructor. Eventually negative rational sequences can be handled by a
sign-normalisation adapter and absolute convergence, but that adapter is not
implemented. Arbitrary oscillating, hypergeometric, or holonomic sequences are
outside this classification. Forge's existing recurrence or telescoping workers
may help prove a majorisation by a rational tail; the classification itself does
not silently extend to those sequences.

\subsection{From summability to asymptotic tail constants: a next extension}

The same residual language suggests a small further development. Once an upper
barrier has established summability, a lower potential satisfying
\[
 B_{\mathrm{low}}(n)-B_{\mathrm{low}}(n+1)\le a_n
\]
can be summed and passed to the limit. For the rational fragment, any positive
constant below $a/(p d_0)$ eventually yields such a lower potential, by the same
leading-coefficient argument with the residual sign reversed. Together with
Theorem~\ref{thm:constructive}, this gives
\[
 \sum_{k=n}^{\infty}\frac{A(k)}{D(k)}
 \sim \frac{a}{p d_0}\,n^{-p}.
\]
Here the argument is a mathematical extension specification, not an implemented
second-polarity checker. The relevant point for Forge is concrete: a single
extra residual polarity could turn a convergence worker into a certified
asymptotic-estimate worker without calling an unchecked symbolic limit engine.
One must first establish summability with the upper barrier; reversing the
residual does not by itself prove that an infinite sum exists.

\section{Synthesising and checking Cauchy moduli}
\label{sec:moduli}

A summability theorem is often less useful to automation than an explicit
answer to ``how far must the finite approximation go?'' A barrier supplies this
answer without searching for a new induction invariant.

Let $\varepsilon\in\Q$ be positive and assume a validated barrier
$C/(n+c)^p$ on $n\ge N$. A simple sufficient choice is
\begin{equation}\label{eq:coarsemodulus}
 K=\max\left\{N,1-c,\left\lfloor\frac C\varepsilon\right\rfloor+1-c\right\}.
\end{equation}
Then $K+c>C/\varepsilon$ and $K+c\ge1$, so
$(K+c)^p\ge K+c>C/\varepsilon$. Therefore $B(K)<\varepsilon$.
For every finite interval beginning at or beyond $K$, the sum of the
nonnegative terms is less than $\varepsilon$. This is an executable Cauchy
modulus for the tail partial sums; the index convention is ``include all terms
with index less than $K$ in the approximation''.

\subsection{The least admitted index}

The simple bound can be wasteful when $p>1$. Let
\[
 q=\left\lfloor C/\varepsilon\right\rfloor,
 \qquad t=\lfloor q^{1/p}\rfloor+1.
\]
Because $t^p$ is an integer, $t^p>C/\varepsilon$ is equivalent to $t^p>q$.
Thus the least index within the certified tail is
\begin{equation}\label{eq:leastmodulus}
 K_* = \max\{N,t-c\}.
\end{equation}
The implementation obtains $t$ by integer doubling and bisection, with no
floating-point root computation.

The witness is checked \emph{without re-running the bisection}. Verify
\[
 K_*\ge N,\qquad C<\varepsilon(K_*+c)^p,
\]
and, unless $K_*=N$, also verify
\[
 C\ge\varepsilon(K_*-1+c)^p.
\]
Strict monotonicity of the positive integer powers establishes leastness among
indices at least $N$. These checks are exact rational inequalities. They do
not ask the checker to trust how the candidate index was discovered.

In the accompanying experiment, 324 rational-accuracy requests were checked
against the 81 barrier records. The least-index routine improved the coarse
index in 115 requests. This measures witness quality inside the same certified
barrier; it is not a claim that the underlying series has no better analytic
tail estimate.

\subsection{Rational accuracy versus an arbitrary real accuracy}

The executable API accepts rational $\varepsilon$. For a Lean goal
$\forall\varepsilon>0,\exists K,\ldots$ with real $\varepsilon$, one may
use an Archimedean or rational-density theorem to obtain a positive rational
accuracy smaller than it, then instantiate the rational witness theorem.
Alternatively the proof may use a noncomputable real-floor expression. Neither
route should be advertised as a numerical algorithm on an arbitrary opaque
real value. This is an instance of Forge's existing executable-versus-logical
context distinction, applied to a new analytic witness~\cite{forge-ledger}.

\section{How the workers cooperate with Forge}
\label{sec:cooperation}

The proposal is not a new general controller. It adds three producers to the
existing obligation mechanism, each with a precise applicability test and a
specific output theorem. Existing source ownership and acceptance contracts
remain prerequisites, not contributions to be rediscovered
~\cite{forge-neighbours,djex-graphs}.

\subsection{Dispatch based on the shape of the obligation}

\begin{longtable}{@{}P{.29\textwidth}P{.29\textwidth}P{.36\textwidth}@{}}
\toprule
Observed source shape & First action & What to export after checking \\
\midrule\endhead
A compact analytic inequality with an exact zero at a rational anchor &
Translate to the anchor and try a valuation-preserving model before flattening the error &
A quantified bound $f(x)\ge\delta(x-a)^m$, including the anchor and the exact interval. \\
A global exponential-polynomial sign on a half-line &
Try a bounded differential ladder, with existing polynomial positivity available for terminals &
A global sign theorem, not a finite-box observation. \\
An infinite sum or requested convergence rate with eventual rational terms &
Prove the eventual representation; construct a barrier or harmonic minorant &
Summability and explicit tail bound, or a theorem of nonsummability; optionally a modulus witness. \\
An ordinary compact numerical bound with no delicate contact &
Use LeanCert or another already verified enclosure library &
The bound returned by the library under the chosen trust profile. \\
A residual is a perfect square or ordinary polynomial problem &
Return to existing algebraic and polynomial workers &
The polynomial theorem, avoiding unnecessary analytic approximation. \\
\bottomrule
\end{longtable}

The prototype anchor is zero. Translation to a general rational anchor,
two-sided intervals, and proofs involving symbolic parameters are controller
extensions, not capabilities demonstrated by its AST. In particular a two-sided
interval requires parity-aware handling of $(x-a)^m$; the positive-half-interval
proof cannot simply be copied to an odd power on the other side.

\subsection{A cross-worker example}

Suppose the source context contains
\[
 0\le x\le\tfrac12,\qquad x-\sin x\le0,
\]
and the target is $g(x)=g(0)$ for an otherwise arbitrary function $g$.
The anchored worker supplies \eqref{eq:sine-output}; ordered arithmetic then
derives $x=0$, and congruence closes the target. The numerical worker needs no
theory of $g$. The equality worker needs no Taylor series. The interface fact
is the original-context inequality, not a new opaque Boolean result.

For a convergence proof, an existing recurrence or summation worker might
establish $0\le u_n\le A(n)/D(n)$ eventually. The new tail barrier supplies a
summable majorant and an explicit tail estimate. Summability comparison then
transfers the result to $u$. This is more useful than requiring the source
sequence itself to be rational, but the comparison theorem and all index guards
must be present. The prototype tests the rational source only; this transfer is
a specified integration, not another executed corpus.

\subsection{Failure diagnostics that change the next attempt}

Each worker should distinguish a failed mathematical precondition from a
resource limit. For anchored models, a failed logarithm guard should create a
positivity/domain obligation or trigger a change of centre; a negative interval
margin should suggest a higher order, a smaller interval, or a stronger
polynomial quotient check. For ladders, a negative intermediate seed rejects a
path, whereas the interior-double-zero obstruction suggests factorisation or
splitting. For rational tails, a failed coefficient test at a requested cutoff
should suggest moving the cutoff or using a stronger polynomial certificate.

These are routing suggestions, not accepted facts. Conversely, a successfully
checked harmonic minorant is a mathematical divergence result in its named
fragment. It must not be collapsed into the same status as ``no upper barrier
found.''

\subsection{A concrete budget policy}

Use separate budgets for source translation, candidate generation, and proof
reconstruction, as the existing Forge architecture requires. Within the new
lanes, order work by measurable quantities: expression-node count, polynomial
degree, denominator bit length, number of rate blocks, residual degree, and
estimated proof-term expansion. For example, the tail worker should reject an
exponent exceeding the admitted checker limit \emph{before} expanding
$(n+c)^p$.

A practical initial routing policy is: try cheap source simplification first;
run a bounded ladder for a recognised global exponential polynomial; use an
anchored model for exact-contact compact goals; otherwise invoke the existing
verified enclosure worker. Infinite-series goals enter the tail lane only
after an eventual source equation is proved. This is a deterministic policy
that can be evaluated against the same workers without communication; a
learned scheduler is neither required nor tested here.

\section{Libraries, version boundaries, and proof reconstruction}
\label{sec:libraries}

\subsection{Reuse LeanCert; do not write a second interval library}

The inspected LeanCert README describes certified interval evaluation,
transcendental functions, automatic differentiation, root certificates,
optimisation, finite sums, integrals, natural-number tail bounds, and checked
unary enclosure extension rules. It exposes, among other names,
\code{LeanCert.evalInterval}, \code{LeanCert.evalInterval_correct},
\code{LeanCert.API.Bounds}, and \code{LeanCert.API.Optimization}
~\cite{leancert-readme}. It also already uses Taylor-model machinery. Therefore
this article does \emph{not} claim that analytic enclosures, tails in every
sense, or Taylor models are absent from the Lean ecosystem.

The recommended division of labour is specific. Use LeanCert's checked
primitive enclosures and ordinary box bounds. Let Forge select an exact-contact
quotient, organise a global differential ladder, or request a series barrier.
Where a library API already proves the exact desired anchored statement, use
that theorem instead of the prototype's weaker primitive majorant. A full
capability comparison against LeanCert was not performed, and no exclusive
capability claim relative to it is justified by these experiments.

The Python package is useful for independent numerical search and advertises a
bundled Lean bridge on supported platforms. Installing it was attempted here,
but network/DNS access in the execution container failed. The bridge was not
run. The executable prototypes in this archive therefore do not depend on it.

\subsection{The observed toolchain mismatch}

Forge's README pins Lean \code{v4.34.0} and Mathlib revision
\begin{center}\small\code{1cf325a0cf67aca2b04d76b5380ff6a9e410aefa}\end{center}
The inspected LeanCert
\code{lean-toolchain} says \code{v4.33.1}, and its Lake dependency selects
Mathlib \code{v4.33.1}~\cite{forge-readme,leancert-build}.
These are not a tested single dependency set. This archive deliberately does
not invent a Lake manifest asserting that they build together.

There are two honest integration options. Port a narrowly selected LeanCert
adapter to Forge's pin and compile it there, or develop the adapter against
LeanCert's pin while clearly postponing the Forge merge. An out-of-process
worker can exchange rational \emph{data} across those versions, but proof
acceptance must occur in the target environment. A proof term from another
version is not accepted merely because its printed theorem statement matches.

\subsection{Two trust choices, and one concrete import hazard}

LeanCert documents a default native checking mode using \code{native_decide},
a kernel mode using \code{decide +kernel} with no fallback, and an automatic
mode that can report a native fallback. Forge must request the intended mode
explicitly; a numerical backend choice such as rational versus dyadic is a
different axis~\cite{leancert-readme,leancert-trust}.

A narrower source observation also matters. The inspected
\path{LeanCert/CertifiedBounds/Li2.lean} intentionally states two numerical
bounds with \code{sorry}; its documentation says matching statements are
proved in a separate heavy CI target. The module is explicitly not a forwarding
module~\cite{leancert-li2}. This does not imply that LeanCert's main checkers
are unproved. It does imply that a Forge adapter must not treat every theorem
under a library's namespace as equally suitable proof authority.

Use the actual verified declaration, or reconstruct the needed result, rather
than importing the lightweight placeholder theorem. Separately proving an equal
statement does not replace the proof body of a declaration that depends on
\code{sorryAx}. The appropriate audit follows the transitive dependencies of
the \emph{accepted theorem}; the mere presence of an unused placeholder elsewhere
in the environment is not itself a dependency of that theorem. No Li2 module is
needed for the algorithms proposed here.

\subsection{A concrete dependency map}

\begin{longtable}{@{}P{.24\textwidth}P{.33\textwidth}P{.36\textwidth}@{}}
\toprule
Component & Reuse & New work actually required \\
\midrule\endhead
Lean core and existing Forge controller &
Original goal identity, expression construction, local hypotheses, acceptance &
Register three applicability predicates and their source-specific residual obligations. \\
Mathlib polynomial and tactic infrastructure &
Exact polynomial identities, ordered arithmetic, powers, finite sums &
A reflected shifted-coefficient positivity theorem and a small certificate decoder. \\
Mathlib calculus/Taylor modules &
Derivative rules, Taylor theory, comparison infrastructure &
Connect an admitted jet or differential ladder to the exact source function. \\
Mathlib infinite-sum infrastructure &
Convergence of nonnegative bounded partial sums, finite-prefix transport, comparison &
A generic barrier theorem, harmonic-minorant theorem, and index-safe instantiation. \\
LeanCert &
Verified elementary enclosures, ordinary compact numerical bounds, optional AD &
A version-pinned adapter preserving the anchor and selecting a declared trust mode. \\
Python \code{fractions}, integer arithmetic &
Portable exact search and prototype replay &
Already implemented; never substitute Python acceptance for the Lean bridge. \\
SymPy and mpmath &
Independent symbolic diagnostics and high-precision numerical samples &
Testing only. They do not appear in the shipped search/replay dependency chain. \\
\bottomrule
\end{longtable}

Mathlib's documented Taylor definitions and derivative-comparison lemmas are
source-level integration targets, not proof scripts compiled in this session
~\cite{mathlib-taylor,mathlib-meanvalue}. Verify exact names and signatures under
the selected pin before implementing an adapter. In particular, assumptions
about differentiability within an interval and endpoint continuity must not be
lost when selecting a convenient comparison theorem.

\section{The missing Lean layer, specified narrowly}
\label{sec:lean}

The easiest first implementation among these \emph{new} lanes is the rational
tail worker. It needs exact polynomial checking, finite telescoping, and a
summability theorem, but no large catalogue of analytic primitive rules. This
recommendation should not delay Forge's already proposed narrow end-to-end
milestone: the existing finite-algebra cover remains a reasonable first
integration test~\cite{forge-conclusion}. The new lanes are modular extensions,
not a reason to restart the project.

\subsection{Proposed module boundaries}

\begin{lstlisting}
Forge/Analysis/Source.lean
Forge/Analysis/Anchor/Semantics.lean
Forge/Analysis/Anchor/Checker.lean
Forge/Analysis/Anchor/Soundness.lean
Forge/Analysis/Differential/Ladder.lean
Forge/Analysis/RationalTail/Shift.lean
Forge/Analysis/RationalTail/Checker.lean
Forge/Analysis/RationalTail/Soundness.lean
Forge/Analysis/RationalTail/Modulus.lean
Forge/Analysis/Adapters/LeanCert.lean
Forge/Analysis/Route.lean
\end{lstlisting}

These are proposed paths, not existing source names. The architecture keeps
source translation distinct from checker soundness. A reflected polynomial
check can be perfectly correct while proving something about the wrong source
sequence or wrong cast.

\subsection{The three theorem families to formalise}

For an anchored enclosure, prove a generic theorem with the semantic premise
\eqref{eq:anchorsem} and the checked quotient bound. Separately prove every AST
constructor preserves the enclosure. The final source theorem combines the
reification equality, the model soundness theorem, and
Theorem~\ref{thm:anchored}. No theorem may assume ``the Python model is
correct'' as an axiom.

For ladders, first prove Lemma~\ref{lem:factorstep}; then prove soundness of a
list of admitted stages by induction on the list. The exact derivative identity
for an exponential polynomial is reconstructed from the coefficient data using
ordinary derivative rules and ring normalisation. The terminal polynomial signs
can be proved by existing polynomial tools, rather than baking the current
coefficient cone into the outer ladder theorem.

For tails, first prove \eqref{eq:finitetail} by induction on the finite range
length. Next prove the shifted-coefficient sign lemma; next connect
\eqref{eq:residual} to the barrier step using positive denominator guards.
Then prove summability and the tail estimate from bounded monotone partial sums.
Only after these steps should the tactic construct an infinite-sum equality or
inequality in the source goal. The harmonic-minorant theorem is a separate
soundness theorem with a separate conclusion.

\subsection{Source obligations that must be explicit}

For analytic models, retain the exact source variable, function occurrence,
rational scale, anchor, interval, and side of the anchor. A symbolic scale is
not a rational constant simply because a current numerical model has assigned
it a value. A logarithm node requires its original positive-argument proof.
Unsupported composition must return an obligation or remain opaque; it must not
be normalised into a different admitted primitive.

For tails, retain the sequence and a proof of the eventual equation, the
starting index, the denominator sign theorem, and the casts from naturals to
reals. Natural subtraction is not field subtraction. A requested bound from
index $N_0$ cannot be replaced by a bound from a larger discovered $N$ without
adding the intervening finite sum. The construction can prove convergence after
moving the cutoff, but a quantitative requested tail bound has stricter indexing
requirements.

For ladders, retain the exact real-valued source function and the half-line.
Boundary values must be exact rational evaluations, not floating-point
approximations. The same differential identity on an interval that misses zero
does not license use of the seed at zero without a bridge over the gap.

These requirements are the analytic application of source-owned evidence
already developed in Djex and Leant. Their source-typed graph design explicitly
associates evidence with a particular candidate and checked inventory, rather
than inferring authority from a displayed global name~\cite{djex-graphs}.
Nothing in this proposal resolves their separate open Church-indexing or
universe-composition synthesis tasks.

\subsection{What the supplied Lean file demonstrates}

The archive contains \code{lean/SemanticBridges.lean}, with two small candidate
lemmas: nonnegativity from an already established anchored factorisation, and
the finite telescoping consequence of an already established barrier step.
They contain no admitted proofs. They are deliberately modest: neither lemma
checks a JSON certificate or reifies a source expression. The file was not
compiled here, and it must not be counted as a kernel-checked artifact.

Providing these candidates is useful for fixing the intended proposition and
index convention. Calling them an implemented tactic would not be useful. The
substantive missing work is the checker-soundness/source-bridge composition,
followed by actual kernel acceptance and axiom-dependency inspection.

\section{Experimental evidence and its limits}
\label{sec:experiments}

\subsection{Environment and reproduction}

The recorded run used Python 3.13.5 on Linux x86-64 with glibc 2.41. The optional
oracle run used SymPy 1.14.0 and mpmath 1.3.0 at 120 decimal digits. The random
seed is 20260915. The exact commands, problem subjects, certificates, mutation
outcomes, and unindented JSON byte measurements are retained. No Lean,
\code{lake}, or \code{elan} executable was found. Attempts to install or clone
through the execution container failed because external host names could not
be resolved; GitHub source inspection itself was available through the connected
read interface.

\begin{longtable}{@{}P{.32\textwidth}r r P{.26\textwidth}@{}}
\toprule
Certificate lane & Records & Distinct subjects & Largest JSON record \\
\midrule\endhead
Anchored enclosures & 62 & 61 & 339 bytes; order~8 \\
Differential ladders & 43 & 42 & 671 bytes; depth~4 \\
Upper tail barriers & 81 & 67 & 878 bytes; cutoff~21 \\
Harmonic minorants & 60 & 60 & 387 bytes; cutoff~12 \\
\midrule
All certificate records & 246 & 230 & Not a Lean proof-size metric \\
\bottomrule
\end{longtable}

The totals in this table are counts of the same object, namely certificate
records, so adding them is meaningful. They are not added to unit tests,
mutations, numeric samples, or finite-sum checks. Nor are they a solved-goal
benchmark. The repeated subjects are disclosed in
\code{results/duplicate-audit.json}.

\subsection{Anchored corpus and ablation}

The corpus consists of 48 scale/radius combinations of four elementary gaps,
eight products of two such gaps, and six higher-order exponential remainders.
The elementary families are
\[
 e^{ax}-1-ax,\quad ax-\sin(ax),\quad ax-\log(1+ax),\quad ax-\arctan(ax),
\]
with $a\in\{1/4,1/2,1,2\}$ and $ab\in\{1/4,1/2,3/4\}$.
The higher-order subfamily repeats one subject from the elementary family.

All 62 records are accepted. The factor-erasing ablation uses the \emph{same}
expression, box, chosen order, and remainder model, but applies the lower bound
from $H_b(p)+[0,b^K]I$ without valuation extraction. It succeeds in 8 of the 62
records. Those eight are the product subfamily: the model already has zero
truncated polynomial and a positive remainder interval, so the simple bound
suffices.

On distinct subjects the figures are 61 accepted by anchored extraction versus
8 accepted by the ablation. This isolates a representation loss, not a
comparison against a well-tuned interval prover. Signed remainder rules,
symbolic factorisation, a stronger Taylor implementation, or an existing
library theorem can also solve some or all of these examples. In particular,
the product examples have intentionally simple mathematical structure. The
ablation is not evidence that the new worker outperforms LeanCert.

SymPy independently reconstructs all 62 Taylor polynomials from the function
expressions. A separate mpmath evaluator checks 620 nonzero sample points against
the stored remainder intervals. Sampling is diagnostic: it cannot certify a
universal interval claim and is not used by the acceptance checker.

\subsection{Ladder corpus and its planted component}

There are seven hand-described exponential-remainder records, including one
repeated subject, and 36 planted examples. A planted example begins with a
terminal-cone exponential polynomial $g$ and repeatedly solves
\[
 (D-\alpha)f=g,\qquad f(0)=c\ge0
\]
exactly within the exponential-polynomial language. This deliberately creates
inputs with known ladder certificates. The search is then run from the resulting
source function, without being given the planted rate list.

All 43 records are accepted. The checker validates the displayed stages, not the
construction history. SymPy cross-checks 68 differential identities in the
returned ladders. Planted examples test search/replay consistency and sensitivity
to degree/rate structure. They are not representative samples from arbitrary
user inequalities. The always-nonnegative function $(e^x-2)^2$ is retained as
an expected \Unknown{} control, and its obstruction is proved in
Section~\ref{sec:ladder} rather than explained away as insufficient fuel.

\subsection{Tail corpus, witnesses, and exact controls}

The rational-function generator creates 120 subjects with random coefficients
and positive leading coefficients. Numerator degrees lie between zero and four;
requested degree offsets cycle through $-2,-1,0,1,2,3,4,5$, with denominator
degree clamped below by zero. The construction produces 60 upper barriers and
60 harmonic minorants. Unlike the ladder corpus, these rational subjects are
not obtained by inverting the certificate equations; the completeness theorem
explains why the appropriate certificate exists.

For the 60 constructed upper barriers, 4,800 exact finite partial-sum
inequalities are checked, and 240 coarse rational-accuracy witnesses are checked.
For the 60 divergence records, 2,400 exact pointwise harmonic comparisons are
checked. These finite checks are sanity checks of the infinite certificate,
not a substitute for its mathematical proof.

An additional 21 barrier records optimise $1/n^s$ for $s=2,\ldots,8$ at
cutoffs $2,5,10$. They concern seven sequences and three quantitative requests
per sequence. The optional SymPy run verifies the shifted residual/denominator
identities in all 141 tail records. Finally 324 least-index witnesses, four for
each of the 81 upper barrier records, pass their exact two-inequality checks;
115 improve the coarse witness.

\subsection{Replay, mutations, and explicit non-successes}

The standalone replay is invoked with Python's \code{-S} option, disabling
site packages. Its import guard rejects attempts to import the search modules
or numerical/symbolic libraries. All 246 records replay. This establishes a
separable implementation path, not a formally verified checker.

Three evidence-bearing mutations are applied to every record: source subject
alteration, a bound/seed/minorant alteration, and a polynomial or stage alteration.
All 738 are rejected. Additional unit tests reject duplicate JSON keys,
floating-point JSON, noncanonical rational strings, Boolean values masquerading
as integer cutoffs, and a denominator vanishing at the requested cutoff.
The standard-library test suite contains 25 test methods; all pass.

The recorded analytic controls include four false elementary inequalities, an
unmet logarithm-series domain, and the true trigonometric identity discussed
earlier. All return \Unknown{} under the stated caps. The global even-root
square is also \Unknown{}. These outcomes are not packaged as refutations.
Only a separately validated harmonic minorant authorises the divergence
conclusion in the rational-series lane.

\subsection{What remains unmeasured}

There is no compiled Lean checker, source reifier, Forge tactic call, or
LeanCert adapter in this experiment. There is no Lean axiom audit, no kernel
replay time, no proof-term size measurement, and no comparison against
\code{grind}, \code{nlinarith}, a tuned LeanCert configuration, or a
noncommunicating portfolio. The small JSON certificates suggest a manageable
transport format; they say nothing by themselves about the size or cost of the
resulting Lean proof terms.

The Python checker and search use different representations but share the same
mathematical design. A common error in the derivation could survive both. The
paper proofs, symbolic oracles, numeric diagnostics, and mutations attack
different failure modes; none makes the checker formally verified. Likewise,
input caps are basic resource guards, not a security audit against adversarial
resource exhaustion.

\section{An implementation and evaluation programme}
\label{sec:programme}

\subsection{Acceptance gates rather than speculative deadlines}

The first gate is the finite barrier lemma and a polynomial sign checker in
Lean under one chosen toolchain. Acceptance requires a source theorem about
$\sum_{i=0}^{m-1}a(N+i)$, not merely a coefficient identity, and a recorded axiom
inventory. Wrong-denominator and off-by-one variants must be rejected.

The second gate adds the generic infinite-tail theorem and finite-prefix
transport. Acceptance requires both a positive convergence example and a
harmonic-divergence example, with the exact original source sequence preserved.
A passing finite-sum lemma alone does not close this gate.

The third gate adds coefficient-driven synthesis and least-index witnesses.
Acceptance requires replay with the search disabled, including verification of
the strict inequality at the chosen index and the predecessor inequality when
leastness is claimed. The convergence theorem must continue to compile after
removing Python and all optional search dependencies.

The fourth gate proves and instantiates the integrating-factor step. A global
exponential remainder should compile from a generated ladder, while a fabricated
ladder for the even-root square must fail. A proof over a finite box is not an
acceptable substitute for the original half-line goal.

The fifth gate connects anchored models to source functions, reusing checked
primitive enclosures where possible. Acceptance requires an endpoint-contact
example, a product with high-order cancellation, a failed logarithm guard, and
an exact identity routed away from interval approximation. Only then should
these workers enter the normal Forge scheduler.

\subsection{The benchmark that would demonstrate practical value}

Evaluate four configurations on the same held-out source goals and identical
premise sets: existing Forge without the new lanes; the same new workers as a
noncommunicating portfolio; the communicating integration; and strong directly
invoked alternatives including tuned LeanCert for supported numeric goals.
Fix compiler versions, imports, theorem-retrieval restrictions, and budgets.
Remove any target theorem or trivially equivalent generated lemma from retrieval.

Stratify the analytic goals by anchor order and number of cancelling terms;
stratify global inequalities by rate blocks and terminal polynomial difficulty;
stratify tails by rational degree gap, mixed coefficients, requested cutoff,
and accuracy. Include false, unsupported, and deliberately difficult true
controls. Report successful kernel proofs, unknowns, incorrect acceptances,
source-translation failures, search time, replay time, proof size, and imported
axioms separately.

The main hypothesis is not that every new worker is stronger than every
existing tactic. It is that preserving contact order, finding a global
comparison chain, and generating quantitative tail witnesses will close source
obligations that currently require manual mathematical restructuring. That
hypothesis becomes credible only when tested on actual Lean goals, with the
source and proof boundary included in the measurement.

\section{Conclusion}

Three additions can broaden Forge without rebuilding its existing foundations.
Anchored remainders retain the local order of vanishing that a uniform numerical
error can erase. Integrating-factor ladders turn a global analytic sign into a
finite chain of exact differential identities. Rational tail barriers turn
infinite convergence into shifted polynomial signs, with explicit divergence
certificates on the complementary eventually positive rational fragment and
checked quantitative witnesses when convergence holds.

The constructions are concrete enough to implement and test: exact ASTs,
coefficient formulas, finite certificate formats, bounded search policies,
checker rules, soundness arguments, and a constructive completeness theorem are
all supplied. Their limitations are equally concrete: a narrow analytic grammar,
loss of correlations, a first-order global-comparison obstruction, rational-only
tail classification, and an unimplemented Lean source/checker bridge.

The recommended contribution to the repository is therefore a small new
\code{Analysis} extension area with the supplied prototypes and a gate-driven
formalisation plan, not a replacement for \code{grind} and not another general
trust-system redesign. The next decisive result is a kernel-checked source
theorem through one of these lanes. The evidence supplied here makes that
implementation specific, while leaving its unfinished status visible.

\appendix
\section{Certificate schemas and replay contract}
\label{app:schemas}

Rationals are canonical strings \code{"numerator/denominator"}, with positive
denominator, reduced form, and no leading plus or redundant zeros. Polynomial
coefficients are in increasing degree order; the zero polynomial is
\code{["0/1"]}. The search uses dense tuples, while replay converts them to
sparse maps with zero coefficients removed. Exponential blocks are sorted by
rational rate and must not repeat a rate.

Every certificate carries a subject for transport, but the verifier also takes
the expected subject \emph{externally}. It first requires exact agreement.
In a Lean integration that external subject must be derived from the original
goal and its proved source translation, not read back from the candidate.

\begin{lstlisting}
anchored:
  kind, subject={expr,b}, order, polynomial,
  error=[lower,upper], valuation, margin

ladder:
  kind, subject={exp_poly,domain="x>=0"},
  rates, stages, seeds, search_nodes

barrier:
  kind, subject={numerator,denominator,
                 interpretation="eventual-rational-tail"},
  cutoff, shift, power, constant, shifted_numerator,
  shifted_denominator, shifted_residual, tail_bound

divergence:
  kind, subject={numerator,denominator,
                 interpretation="eventual-rational-tail"},
  cutoff, minorant, shifted_denominator, shifted_residual
\end{lstlisting}

The ladder's \code{search_nodes} is descriptive metadata. It is bounded by the
decoder but is not evidence of correctness or completeness. Every mathematical
field is checked or recomputed. Unexpected fields are rejected by the current
version, so adding a new certificate rule requires an explicit schema change.

Principal replay limits are a 2 MB JSON envelope, 400 characters per rational,
1,024 bits in each admitted rational numerator/denominator, polynomial input
length at most 193, anchored order at most 40, at most 32 exponential blocks and
32 ladder steps, tail exponent at most 16, and tail cutoff at most $10^6$.
These caps bound admitted inputs, not the bit lengths of every intermediate
arithmetic result. A production checker should add independently accounted
intermediate-allocation and proof-reconstruction limits.

\section{Using the archive}
\label{app:use}

From the archive root, no additional Python package is needed for these commands:
\begin{lstlisting}
python -m unittest discover -s tests -v
python -S tools/verify.py
python tools/experiments.py --output reproduced-results
python tools/modulus_checks.py --output reproduced-moduli.json
\end{lstlisting}
The experiment runner refuses an existing output directory, preventing a fresh
run from silently replacing the recorded evidence. It writes a failure receipt
if an assertion fails. The optional oracle diagnostics require the versions in
\code{requirements-testing.txt}:
\begin{lstlisting}
python -m pip install -r requirements-testing.txt
python tools/oracle_checks.py --output reproduced-oracle.json
\end{lstlisting}

\Needspace{15\baselineskip}
A minimal programmatic use is:
\begin{lstlisting}
from fractions import Fraction as Q
from forge_analytic import jets as J
from forge_analytic.checker import verify

expr = J.minus(J.P(0, 1), J.atom('sin'))
expected = {'expr': expr, 'b': '1/2'}
candidate = J.search(expr, Q(1, 2))
if candidate is not None and verify(expected, candidate):
    print(candidate['valuation'], candidate['margin'])
# 3 31/216
# This is Python evidence, not a Lean proof.
\end{lstlisting}
Set \code{PYTHONPATH=prototype} or insert that directory in the Python module
path before using the package directly. The command-line tools do this
internally.

To rebuild the article:
\begin{lstlisting}
cd article
pdflatex -interaction=nonstopmode -halt-on-error forge_analytic_extensions.tex
pdflatex -interaction=nonstopmode -halt-on-error forge_analytic_extensions.tex
pdflatex -interaction=nonstopmode -halt-on-error forge_analytic_extensions.tex
\end{lstlisting}
The document uses ordinary pdf\LaTeX{} packages, no shell escape, no external
bibliography processor, and no redistributed font files.

\section{Source audit and attribution}

The source inventory in \code{results/source-audit.json} records the repository
revision and paths inspected. The audit is scoped to the items discussed in
Section~\ref{sec:delta}; it is not an exhaustive textual similarity audit of all
frozen proposals. The numerical literature is credited for the underlying
techniques. The soundness arguments and specialisations in this article are
written out so the reader can assess the exact admitted fragment without
relying on an informal appeal to a solver.

The code in this archive is original prototype code for these extensions.
No third-party repository archive, package binary, or font is redistributed.
Referenced projects and papers retain their own licences. The accompanying
\code{STATUS.md} distinguishes executed Python artifacts, optional oracle
checks, uncompiled Lean candidates, and unimplemented integration work.

\begin{thebibliography}{99}
\bibitem{forge-readme}
V. Reshetnikov, \emph{Forge}, README at audited revision
\code{c98e47c5} (full identifier in the source audit), inspected 15 September 2026.
\url{https://github.com/VladimirReshetnikov/Forge/blob/c98e47c5f804e92880fc1d0e378c1b95832b685c/README.md}.

\bibitem{forge-ledger}
Forge, \emph{Provenance: what came from where}, merged article appendix B,
same revision. \url{https://github.com/VladimirReshetnikov/Forge/blob/c98e47c5f804e92880fc1d0e378c1b95832b685c/article/sections/B-provenance.tex}.

\bibitem{forge-conclusion}
Forge, \emph{Assessment and recommended starting point}, merged article
section 12, same revision.
\url{https://github.com/VladimirReshetnikov/Forge/blob/c98e47c5f804e92880fc1d0e378c1b95832b685c/article/sections/12-conclusion.tex}.

\bibitem{forge-baseline}
Forge, \emph{A strong baseline, not a straw man}, merged article section 2,
same revision.
\url{https://github.com/VladimirReshetnikov/Forge/blob/c98e47c5f804e92880fc1d0e378c1b95832b685c/article/sections/02-baseline.tex}.

\bibitem{forge-nonlinear}
Forge, merged article section 7, nonlinear certificates, especially the
univariate, product-envelope, and domain discussion, same revision.
\url{https://github.com/VladimirReshetnikov/Forge/blob/c98e47c5f804e92880fc1d0e378c1b95832b685c/article/sections/07-nonlinear.tex}.

\bibitem{forge-neighbours}
Forge, \emph{Ideas from Leant and Djex}, same revision.
\url{https://github.com/VladimirReshetnikov/Forge/blob/c98e47c5f804e92880fc1d0e378c1b95832b685c/docs/IDEAS-FROM-LEANT-DJEX.md}.

\bibitem{djex-graphs}
V. Reshetnikov, Djex, \emph{Djinn source-typed evidence graph}, inspected
15 September 2026; source blob
\code{7c1dca61d49f6e113c44fb9aba73d3203296cd5d}.
\url{https://github.com/VladimirReshetnikov/Djex/blob/main/docs/source-typed-evidence-graph.md}.

\bibitem{leant}
V. Reshetnikov, \emph{Leant---a Djex-based synthesis REPL for Lean 4},
README, inspected 15 September 2026.
\url{https://github.com/VladimirReshetnikov/Leant}.

\bibitem{autobound}
M. Streeter and J. V. Dillon, \emph{Automatically Bounding the Taylor Remainder
Series: Tighter Bounds and New Applications}, arXiv:2212.11429, version 3,
2 August 2023. \url{https://arxiv.org/abs/2212.11429v3}.

\bibitem{verified-intervals}
M. Daumas, D. Lester, and C. Mu\~noz, \emph{Verified Real Number Calculations:
A Library for Interval Arithmetic}, arXiv:0708.3721, 2007.
\url{https://arxiv.org/abs/0708.3721}.

\bibitem{leancert-readme}
LeanCert contributors, \emph{LeanCert: Certified numerics for Lean 4}, README,
inspected 15 September 2026; source blob
\code{8a65c45afc1763a36d5d406692779daff401b92f}.
\url{https://github.com/alerad/leancert/blob/main/README.md}.

\bibitem{leancert-build}
LeanCert contributors, \code{lean-toolchain} and \code{lakefile.toml}, inspected
15 September 2026; Lean and Mathlib pins \code{v4.33.1}.
\url{https://github.com/alerad/leancert/blob/main/lean-toolchain};
\url{https://github.com/alerad/leancert/blob/main/lakefile.toml}.

\bibitem{leancert-trust}
LeanCert, \emph{Trust model}, official documentation, inspected
15 September 2026.
\url{https://docs.leancert.io/architecture/trust-model/}.

\bibitem{leancert-li2}
LeanCert contributors, \code{LeanCert/CertifiedBounds/Li2.lean}, inspected
15 September 2026; source blob
\code{5833f2515cbb42d1a2bd21fcdaedc5362be7cbaf}.
\url{https://github.com/alerad/leancert/blob/main/LeanCert/CertifiedBounds/Li2.lean}.

\bibitem{mathlib-taylor}
Mathlib contributors, \code{Mathlib.Analysis.Calculus.Taylor}, official API
reference, inspected 15 September 2026.
\url{https://leanprover-community.github.io/mathlib4_docs/Mathlib/Analysis/Calculus/Taylor.html}.

\bibitem{mathlib-meanvalue}
Mathlib contributors, \code{Mathlib.Analysis.Calculus.MeanValue}, official API
reference, especially derivative comparison on finite intervals, inspected
15 September 2026.
\url{https://leanprover-community.github.io/mathlib4_docs/Mathlib/Analysis/Calculus/MeanValue.html}.

\bibitem{mathlib-infinitesum}
Mathlib contributors, \code{Mathlib.Topology.Algebra.InfiniteSum.Order},
official API reference, inspected 15 September 2026.
\url{https://leanprover-community.github.io/mathlib4_docs/Mathlib/Topology/Algebra/InfiniteSum/Order.html}.

\bibitem{mathlib-sumintegral}
Mathlib contributors, \code{Mathlib.Analysis.SumIntegralComparisons}, official
API reference, inspected 15 September 2026.
\url{https://leanprover-community.github.io/mathlib4_docs/Mathlib/Analysis/SumIntegralComparisons.html}.

\bibitem{grind}
Lean language reference, \emph{The grind tactic}, current official reference,
inspected 15 September 2026.
\url{https://lean-lang.org/doc/reference/latest/The--grind--tactic/}.
\end{thebibliography}
\end{document}
